\documentclass[12pt,a4paper]{article}

\usepackage[utf8]{inputenc}
\usepackage[T1]{fontenc}
\usepackage[polish]{babel}
\usepackage{geometry}
\geometry{lmargin=25mm, rmargin=25mm, tmargin=25mm, bmargin=25mm}
\usepackage{setspace}
\onehalfspacing
\usepackage{hyperref}
\usepackage{enumitem}

% Wpis: autor → tytuł → opis wykorzystania (osobne linie).
\newcommand{\lit}[3]{%
  \item
  \textbf{#1}\\[0.15em]
  \textit{#2}\\[0.35em]
  #3%
}

\begin{document}

\begin{center}
    {\Large \textbf{Zestawienie literatury i plan jej
    wykorzystania w pracy}}\\
    \vspace{0.3cm}
    {\large Rafał Bazan (nr albumu: 176768)}\\
\end{center}

\vspace{0.5cm}

\section*{1. Bariery kognitywne i interakcja człowiek-rój (HSI)}
\textit{Literatura wykorzystana w Rozdziale 1. do
uzasadnienia potrzeby zmiany sposobu sterowania rojem
z ręcznego na wspomagany sztuczną inteligencją.}

\begin{itemize}[leftmargin=*, itemsep=1.1em]
    \lit{A. Kolling et al. (2016)}
        {Human interaction with robot swarms: A survey}
        {Praca definiuje wymóg przejścia na sterowanie
        grupowe o stałej złożoności $O(1)$. Służy jako
        podstawa do uargumentowania, dlaczego operator
        powinien wydawać ogólne dyrektywy całemu rojowi,
        zamiast kontrolować każdą jednostkę z osobna.}

    \lit{J. Y. C. Chen (2010) oraz
        Defense Intelligence Agency (2010)}
        {Supervisory Control of Unmanned Vehicles /
        Cognitive Limits on Simultaneous Control of
        Multiple Unmanned Spacecraft}
        {Raporty wojskowe pokazujące limity ludzkiej
        uwagi przy zarządzaniu wieloma obiektami naraz.
        Stanowią uzasadnienie, dlaczego interfejs oparty
        na modelach językowych jest optymalnym
        rozwiązaniem zapobiegającym przeciążeniu
        operatora.}

    \lit{C. Smith et al. (2021) / L. Jia et al. (2026)}
        {A Cognitive Walkthrough of Multiple Drone
        Delivery Operations / Understanding Situation
        Awareness in Multi-UAV Supervision}
        {Artykuły analizujące poziomy świadomości
        sytuacyjnej i zmęczenia operatorów floty
        bezzałogowców. Mogą posłużyć w dalszych etapach
        jako punkt odniesienia przy wyznaczaniu metryk
        oceny obciążenia człowieka podczas testów.}
\end{itemize}

\section*{2. Embodied AI i modele językowe (LLM)}
\textit{Podstawa warstwy decyzyjnej projektu
(Rozdział 3). Publikacje tłumaczą, jak semantycznie
powiązać mowę z systemem operacyjnym robota.}

\begin{itemize}[leftmargin=*, itemsep=1.1em]
    \lit{K. Rachwał et al. (Robotec.AI, 2025)}
        {RAI: Flexible Agent Framework for Embodied AI}
        {Dokumentacja frameworku RAI do integracji
        modeli generatywnych z architekturą ROS2.
        Stanowi punkt wyjścia do zaplanowania warstwy
        pośredniczącej między modelami językowymi
        a komponentami wykonawczymi systemu.}

    \lit{C. Sun et al. (2025)}
        {LLM-Based Multi-Agent Decision-Making:
        Challenges and Future Directions}
        {Praca o wykorzystaniu dużych modeli językowych
        w podejmowaniu decyzji przez grupy robotów.
        Pomaga w strukturyzacji podpowiedzi (promptów)
        i logiki decyzyjnej ukierunkowanej na bezpieczny
        dobór strategii nawigacyjnej.}

    \lit{M. Pomianek (2026)}
        {Analiza literatury HRI}
        {Przegląd zaawansowanych interfejsów
        komunikacyjnych. Uzasadnia zastosowanie koncepcji
        Edge Computing (przetwarzania lokalnego) pod
        kątem minimalizacji opóźnień w przesyłaniu
        poleceń i ochrony prywatności danych
        biometrycznych.}
\end{itemize}

\section*{3. Inteligencja stadna i algorytmy nawigacyjne}
\textit{Literatura opisująca mechanikę ruchu grupy
dronów (część matematyczna Rozdziału 1 i implementacja
w Rozdziale 3).}

\begin{itemize}[leftmargin=*, itemsep=1.1em]
    \lit{Zhao et al. (2018)}
        {Autonomous Decision Making for UAV Cooperative
        Pursuit-Evasion Game with Reinforcement Learning}
        {Artykuł o sterowaniu rojem przy użyciu uczenia
        ze wzmocnieniem. Służy jako punkt odniesienia
        pokazujący zalety podejścia hybrydowego
        (połączenie warstwy LLM z deterministycznymi
        algorytmami) pod kątem przewidywalności systemu
        w sytuacjach krytycznych.}

    \lit{Ł. Strąk (2017)}
        {Adaptacyjny algorytm optymalizacji stadnej
        cząsteczek dla dynamicznego problemu
        komiwojażera}
        {Rozprawa doktorska wprowadzająca pojęcie
        „pamięci feromonowej”. Publikacja ta stanowi
        cenną inspirację teoretyczną przy projektowaniu
        mechanizmów zapobiegania utykaniu roju
        w minimach lokalnych i ślepych zaułkach mapy.}

    \lit{E. Bonabeau (1999) / J. Kennedy (2001)}
        {Swarm Intelligence}
        {Klasyczne pozycje opisujące rozproszoną
        samoorganizację w przyrodzie.
        Dostarczają podstaw teoretycznych do analizy
        zachowań kolektywnych i odporności rozproszonej
        struktury stada na utratę części jednostek.}
\end{itemize}

\section*{4. Technologie i środowisko symulacyjne
(ROS2 / Nav2)}
\textit{Materiały inżynierskie, które posłużą jako
wytyczne do konfiguracji środowiska badawczego
w Rozdziałach 2. i 3.}

\begin{itemize}[leftmargin=*, itemsep=1.1em]
    \lit{S. Macenski (2020) i A. Bonci (2023)}
        {The Marathon 2: A Navigation System /
        ROS2-Based Frameworks for Increasing Robot
        Autonomy}
        {Opisy stosu nawigacyjnego Navigation2 (Nav2)
        dla autonomicznych robotów mobilnych.
        Przydatne przy analizie Drzew Zachowań
        (Behavior Trees) oraz mechanizmów
        asynchronicznego omijania przeszkód.}

    \lit{Dokumentacja techniczna platformy}
        {DynamicSwarms: ds-crazyflies Documentation}
        {Zbiór specyfikacji technicznych dla
        symulatora dronów Crazyflie. Pozwala na
        uwzględnienie realnych ograniczeń
        kinematycznych i parametrów fizycznych
        platform latających w środowisku testowym.}
\end{itemize}

\end{document}
